\documentclass[10pt,a4paper]{amsart}
\usepackage[margin=19mm]{geometry}
\usepackage{amsmath,amssymb,mathtools}
\usepackage[scr=rsfs,cal=euler]{mathalfa}
\usepackage{xfrac}
\usepackage[unicode,hidelinks,bookmarks=false]{hyperref}
\newcommand{\ie}{i.e.\ }
\newcommand{\cf}{cf.\ }
\newcommand{\resp}{respectively\ }
\newcommand{\Subsection}{\subsection}
\providecommand{\nobreakdash}{\nobreak\hbox{-}}
\setlength{\emergencystretch}{2em}
\multlinegap0pt
\usepackage{bm}
\usepackage{bbm}
\usepackage{dsfont}
\usepackage{xspace}
\usepackage{cancel}
\usepackage{mathtools}
\usepackage{amsthm}
\makeatletter
\@ifundefined{question}{\newtheorem{question}{Question}}{}
\@ifundefined{proposition}{\newtheorem{proposition}[question]{Proposition}}{}
\makeatother
\newcommand{\abs}[1]{\lvert#1\rvert}
\title[Proving it is impossible; on Erd\H{o}s problem $\# 278$]{Proving it is impossible; on Erd\H{o}s problem $\# 278$}
\author{Stijn Cambie}
\date{Source: arXiv:2508.18270v2 (CC BY 4.0)}
\begin{document}
\maketitle

\noindent\emph{Source and licence.} Proving it is impossible; on Erd\H{o}s problem $\# 278$. Version arXiv:2508.18270v2 (CC BY 4.0), distributed under the Creative Commons Attribution 4.0 International licence, as recorded in the arXiv licence metadata for this exact version and shown on the abstract page https://arxiv.org/abs/2508.18270v2. Full rights notice: licenses/arxiv-2508.18270-CC-BY-4.0.txt.\par\smallskip

\noindent\textit{Editorial scope.} Counted unit: the Proposition whose source label is \texttt{prop:chvatal} in the article (source0.tex lines 259--271, source label \texttt{prop:chvatal}), delivered together with the article's own complete original proof of it (source0.tex lines 273--313, one \texttt{proof} environment of 41 lines). No mathematics is written, repaired or re-derived: statement and proof are the article's own text line for line. Required context retained uncounted: none. Every definition, display and label of those lines is retained unchanged. Omitted from this package and not redistributed: 1--258, 272--272, 314--381. Nothing inside a retained excerpt is omitted or altered: every retained body is delivered exactly as the article's lines are written, so no omission receipt of any kind accompanies this package. Citations: the retained excerpts carry no citation command at all, so no bibliography entry of the article is transcribed and no reference is redistributed. Reference adaptation: none was needed and none was made; every reference inside the delivered unit resolves inside it, so no unresolved reference or citation is produced. Printed slips kept literal: no printed word, symbol, hypothesis or number of the counted statement or of its proof was altered. Layout and class: the article's own class is \texttt{article} with packages including inputenc, fontenc, babel, amsmath, amssymb, amsthm, mathtools, lmodern, geometry, microtype, hyperref, cleveref; the wrapper is the lane's pinned \texttt{amsart} editorial wrapper at 10pt on A4, which restates the theorem-like environments the retained text needs and adds the article's own macro definitions that the retained text needs (\texttt{\textbackslash abs}), with extra wrapper packages (bm, bbm, dsfont, xspace, cancel, mathtools). The delivered numbering is the unit's own. Assets: no retained span contains a figure environment, a graphics-inclusion command, a PDF-page inclusion, a table or a TikZ picture; no image file is copied into this package, the package declares no asset dependency and no image file is redistributed with it. Licence: version arXiv:2508.18270v2 of this article is distributed under the Creative Commons Attribution 4.0 International licence, as recorded in the arXiv licence metadata for this exact version and shown on the abstract page https://arxiv.org/abs/2508.18270v2. A full rights notice is delivered at licenses/arxiv-2508.18270-CC-BY-4.0.txt. Characters: every retained excerpt is pure ASCII, so the delivered engine, XeLaTeX with its default Latin Modern text font, reports no missing character.
\begin{proposition}\label{prop:chvatal}
The weights $(w_i)$ and the integers $(A_i)$ induce the same strict ordering
of all subset sums.  Moreover, the integers $(A_i)$ satisfy the following
properties:
\begin{enumerate}
 \item $\sum_{i\in I}A_i<\frac12\sum_{i=1}^n A_i$ whenever
       $10\abs I\leq n$;
 \item no integer greater than $1$ divides more than $\lceil n/2\rceil$ of
       the $A_i$;
 \item distinct subsets have distinct sums;
 \item no subset has sum $\left\lfloor\frac12\sum_{i=1}^nA_i\right\rfloor$.
\end{enumerate}
\end{proposition}

\begin{proof}
Let $I\neq J$ and cancel their intersection.  With
$\sigma_i=1$ on $I$, $\sigma_i=-1$ on $J$, and $\sigma_i=0$ otherwise, put
$D=\sum_i\sigma_iw_i$.  Then $\abs{xD}>4n$, while
\[
 \sum_i\sigma_iA_i
 =c\sum_i\sigma_i\lfloor xw_i\rfloor
   +\sum_{i>m}\sigma_i.
\]
The difference between the right-hand side and $cxD$ has absolute value less
than $cn+n$.  Hence it has the same sign as $D$, and its absolute value is
greater than $2n$.  This proves the asserted order equivalence and property~3.

Put $y=xw_n>4n$.  The choice of the primes gives
\[
 \frac{\max_i A_i}{\min_i A_i}
 \leq\frac{cxw_1+1}{c(xw_n-1)}
 <\frac{3y/2+1}{y-1}<3.
\]
If $10\abs I\leq n$, then
\[
 \sum_{i\in I}A_i
 \leq\abs I\max_iA_i
 <3\abs I\min_iA_i
 <(n-\abs I)\min_iA_i
 \leq\sum_{i\notin I}A_i,
\]
which is property~1.

For property~2 it is enough to consider a prime divisor $q$.  If $q$ divides
some $A_i$ with $i\leq m$, then, since
$\lfloor xw_i\rfloor$ divides $c$, every prime factor of $A_i$ divides $c$.
But $A_j\equiv1\pmod c$ for $j>m$, so $q$ divides none of the latter
integers.  A prime dividing only integers in the second group divides at most
$\lceil n/2\rceil$ of them as well.

Finally, if the total sum is even, property~4 follows from property~3 by
comparing a subset with its complement.  If the total is odd, a subset with
the stated sum would differ from its complement by $1$, contradicting the
lower bound $2n$ above.
\end{proof}

\end{document}
